\section{Series de números complejos}
\subsection{Series de números complejos}

\begin{definición} [Serie de números complejos]
    Una serie de números complejos $\sum_{n=0}^{\infty} z_n$ es convergente si la sucesión de sus sumas parciales 
    \[
        (S_n)_{n\in\mathbb{N}} = \left(\sum_{k=0}^{n} z_k\right)_{n\in\mathbb{N}}
    \]
    es convergente en $\mathbb{C}$. En este caso, a $\lim_{n\to\infty} S_n$ se le llama la suma de la serie y se denota por
    \[
        \sum_{n=0}^{\infty} z_n = \lim_{n\to\infty} S_n.
    \]    
\end{definición}

Si la serie $\sum_{n=0}^{\infty} z_n$ no es convergente, se dice que es divergente o que diverge.

Cuando $(S_n)_{n\in\mathbb{N}}$ es convergente, entonces $(\Re(S_n))_{n\in\mathbb{N}}$ y $(\Im(S_n))_{n\in\mathbb{N}}$ son convergentes en $\mathbb{R}$ y se cumple que

\[
    \sum_{n=0}^{\infty} z_n = \sum_{n=0}^{\infty} \Re(z_n) + i \sum_{n=0}^{\infty} \Im(z_n).
\]

tenemos que 

\[
    \sum_{n=1}^{\infty} z_n \text{ converge } \iff \sum_{n=0}^{\infty} \Re(z_n) \text{ y } \sum_{n=0}^{\infty} \Im(z_n) \text{ convergen.}
\]

Es inmediato comprobar que 
\begin{proposición}
    \vspace{-2em}
    \begin{itemize}
        \item Si $\sum_{n=0}^{\infty} z_n$ converge, entonces $z_n \to 0$.
        \item Si $\sum_{n=0}^{\infty} z_n$ converge y $\sum_{n=0}^{\infty} w_n$ converge, entonces $\sum_{n=0}^{\infty} (z_n + w_n)$ converge y
        \[
            \sum_{n=0}^{\infty} (z_n + w_n) = \sum_{n=0}^{\infty} z_n + \sum_{n=0}^{\infty} w_n.
        \]
        \item Si $\sum_{n=0}^{\infty} z_n$ converge, y $c \in \mathbb{C}$, entonces $\sum_{n=0}^{\infty} c z_n$ converge y
        \[
            \sum_{n=0}^{\infty} c z_n = c \sum_{n=0}^{\infty} z_n.
        \]
    \end{itemize}    
\end{proposición}

\begin{definición} [Convergencia absoluta]
    Una serie $\sum_{n=0}^{\infty} z_n$ se dice absolutamente convergente si la serie $\sum_{n=0}^{\infty} |z_n|$ es convergente.
\end{definición}

Se tiene que 

\[
    \sum_{n=0}^{\infty} z_n \text{ es absolutamente convergente } \implies \sum_{n=0}^{\infty} z_n \text{ es convergente.}
\]

el recíproco no siempre es cierto, por ejemplo, la serie armónica alternada $\sum_{n=1}^{\infty} \frac{(-1)^{n+1}}{n}$ es convergente pero no absolutamente convergente.

Se prueba análogamente a como se hace en $\mathbb{R}$ que

\begin{itemize}
    \item Si $\sum_{n=0}^{\infty} z_n$ es absolutamente convergente $\implies$ todas las permutaciones de $\sum_{n=0}^{\infty} z_n$ son absolutamente convergentes y tienen la misma suma.
    \item Si $\sum_{n=0}^{\infty} z_n$ es absolutamente convergente $\implies$ todas sus subseries son absolutamente convergentes.
\end{itemize}

\ejemplo{
    ¿Dónde converge la serie $\sum_{n=0}^{\infty} z^n$?\\
    \begin{itemize}
        \item Si $|z| > 1$, entonces $|z^n| = |z|^n \xrightarrow[n\to\infty]{} \infty$, luego $z_n \not\to 0$ y la serie diverge.
        \item Si $|z| = 1$, entonces $|z^n| = 1$, luego $z_n \not\to 0$ y la serie diverge.
        \item Si $|z| < 1$, entonces como
        \[
            (1-z)(1 + z + z^2 + \ldots + z^n) = 1 - z^{n+1} \implies 1 + z + z^2 + \ldots + z^n = \frac{1 - z^{n+1}}{1-z},
        \]
        Por tanto,
        \[
            \sum_{n=0}^{\infty} z^n = \lim_{N\to\infty} \left(\sum_{n=0}^{N} z^n\right) = \lim_{N\to\infty} \frac{1 - z^{N+1}}{1-z} = \frac{1}{1-z}.
        \]
        Con lo cual $\sum_{n=0}^{\infty} z^n$ converge $\forall D(0,1)$ y diverge en $\mathbb{C} \setminus D(0,1) = \{z \in \mathbb{C} : |z| \geq 1\}$.\\
        Además, $\sum_{n=0}^{\infty} z^n = \frac{1}{1-z} \quad \forall z \in D(0,1)$.
    \end{itemize}
}

\subsection{Series de potencias}
\begin{definición} [Serie de potencias]
    Si $z_0 \in \mathbb{C}$ y $(a_n)_{n\in\mathbb{N}} \subset \mathbb{C}$, diremos que $\sum_{n=0}^{\infty} a_n (z - z_0)^n$ es una \textbf{serie de potencias} con centro en $z_0$.
\end{definición}

\ejemplo{
    Hemos visto que la serie $\sum_{n=0}^{\infty} z^n$ converge en $D(0,1)$. Por otro lado, la serie $\sum_{n=0}^{\infty} n! z^n$ está centrado en $0$ y converge en $z_0 = 0$, pero en ningún otro punto, ya que si $z \neq 0$,
    \[
        |n! z^n| = n! |z|^n \xrightarrow[n\to\infty]{} \infty,
    \]
    pues 
    \[
        \frac{1}{n! |z|^n} \xrightarrow[n\to\infty]{} 0.
    \]
    ya que 
    \[
        \frac{\frac{1}{(n+1)! |z|^{n+1}}}{\frac{1}{n! |z|^n}} = \frac{1}{(n+1)|z|} \xrightarrow[n\to\infty]{} 0 < 1.
    \]
    Sin embargo, la serie $\sum_{n=0}^{\infty} \frac{z^n}{n!}$ converge $\forall z \in \mathbb{C}$, pues converge absolutamente por el \textit{criterio del cociente}:
    \[
        \frac{\frac{|z|^{n+1}}{(n+1)!}}{\frac{|z|^n}{n!}} = \frac{|z|}{n+1} \xrightarrow[n\to\infty]{} 0 < 1.
    \]
}

\begin{proposición}
    Si la serie de potencias $\sum_{n=0}^{\infty} a_n (z - z_0)^n$ converge en $z_1 \neq z_0$, entonces converge absolutamente $\forall z \in D(z_0, |z_1 - z_0|)$.
\end{proposición}

\begin{proof}
    Dado que $\sum_{n=0}^{\infty} a_n (z_1 - z_0)^n$ converge, entonces $a_n (z_1 - z_0)^n \xrightarrow{n\to\infty} 0$, y por tanto, la sucesión $(a_n (z_1 - z_0)^n)_{n\in\mathbb{N}}$ está acotada.\\
    Llamemos $M > 0$ a una cota superior de dicha sucesión. Así, 
    \[
        |a_n(z - z_0)^n| \leq M \quad \forall n \in \mathbb{N} \cup \{0\}.
    \]
    Tomamos $z \in D(z_0, |z_1 - z_0|)$, entonces tenemos que $|z - z_0| < |z_1 - z_0|$ y además, puesto que $z_1 \neq z_0$, se cumple que
    \[
        \left|\frac{z - z_0}{z_1 - z_0}\right| < 1.
    \]
    llamemos $r = \left|\frac{z - z_0}{z_1 - z_0}\right|$. Entonces $0 < r < 1$. Para cada $n \in \mathbb{N} \cup \{0\}$, se cumple
    \[
        0 \leq |a_n (z - z_0)^n| = |a_n| \left|\frac{z - z_0}{z_1 - z_0}\right|^n |z_1 - z_0|^n \leq M \left|\frac{z - z_0}{z_1 - z_0}\right|^n = M r^n.
    \]
    Como $0 < r < 1$, entonces $\sum_{n=0}^{\infty} r^n$ converge $\implies$ $\sum_{n=0}^{\infty} M r^n$ converge. Por el criterio de comparación, se tiene que $\sum_{n=0}^{\infty} |a_n (z - z_0)^n|$ converge, luego $\sum_{n=0}^{\infty} a_n (z - z_0)^n$ converge absolutamente. 
\end{proof}

\textbf{Nota:} Observemos que para la demostración anterior, si $(a_n (z_1 - z_0)^n)_{n\in\mathbb{N}}$ es acotada entonces $\sum_{n=0}^{\infty} a_n (z - z_0)^n$ converge absolutamente $\forall z \in D(z_0, |z_1 - z_0|)$.

\begin{corolario}
    Si la serie de potencias $\sum_{n=0}^{\infty} a_n (z - z_0)^n$ diverge en un punto $z_2 \in \mathbb{C}$, $z_2 \neq z_0$, entonces diverge $\forall z \in \mathbb{C} \setminus D(z_0, |z_2 - z_0|)$.
\end{corolario}

\begin{proof}
    Si $z \in \mathbb{C} \setminus D(z_0, |z_2 - z_0|)$, entonces la sucesión $(a_n (z - z_0)^n)_{n\in\mathbb{N}}$ no está acotada, pues si lo fuera, por la nota anterior, tendríamos que $\sum_{n=0}^{\infty} a_n (z_2 - z_0)^n$ converge absolutamente en $D(z_0, |z - z_0|)$, lo cual es una contradicción pues $z_2 \in D(z_0, |z - z_0|)$.
\end{proof}

\begin{definición}[Radio de convergencia]
    Dada la serie de potencias $\sum_{n=0}^{\infty} a_n (z - z_0)^n$, se define el \textbf{radio de convergencia} como:
    \[
        R = \sup\left\{|z-z_0| : \sum_{n=0}^{\infty} a_n (z - z_0)^n \text{ converge } \forall z \in D(z_0, r)\right\}.
    \]
    entendiendo que $R = \infty$ si el conjunto anterior no está acotado.
\end{definición}

\begin{teorema}
    Sea $\sum_{n=0}^{\infty} a_n (z - z_0)^n$ una serie de potencias. Entonces existe $R \in [0, \infty]$ tal que:
    \begin{enumerate}
        \item $\sum_{n=0}^{\infty} a_n (z - z_0)^n$ converge absolutamente $\forall z \in D(z_0, R)$.
        \item $\sum_{n=0}^{\infty} a_n (z - z_0)^n$ diverge (más aún, la sucesión $(a_n (z - z_0)^n)_{n\in\mathbb{N}}$ no está acotada) $\forall z \in \mathbb{C} \setminus D(z_0, R)$.
    \end{enumerate}
\end{teorema}

Al disco $D(z_0, R)$ se le llama el \textbf{disco de convergencia} de la serie de potencias. En la circunferencia $C(z_0, R) = \{z \in \mathbb{C} : |z - z_0| = R\}$ puede ocurrir cualquier cosa: la serie puede converger en ningún punto, en algunos puntos o en todos los puntos de la circunferencia.

\ejemplo{
    \begin{itemize}
        \item La serie de potencias $\sum_{n=0}^{\infty} z^n$ tiene radio de convergencia $R = 1$ y diverge en todos los puntos de la circunferencia unidad. Pues $|z| = 1 \implies |z^n| = 1 \not\to 0$.
        \item La serie de potencias $\sum_{n=0}^{\infty} \frac{z^n}{n}$ tiene radio de convergencia $R = 1$. Converge en $-1$, pero diverge en $1$.
        \item La serie de potencias $\sum_{n=0}^{\infty} \frac{z^n}{n^2}$ tiene radio de convergencia $R = 1$ y converge en todos los puntos de la circunferencia unidad, pues $\sum_{n=1}^{\infty} \frac{1}{n^2}$ converge. 
    \end{itemize}
}

\begin{teorema}
    Dada la serie de potencias $\sum_{n=0}^{\infty} a_n (z - z_0)^n$ con radio de convergencia $R > 0$, y sea $f$ su función suma, es decir,
    \[
        f(z) = \sum_{n=0}^{\infty} a_n (z - z_0)^n \quad \forall z \in D(z_0, R).
    \]
    Entonces el radio de convergencia de la serie de las derivadas,
    \[
        \sum_{n=1}^{\infty} n a_n (z - z_0)^{n-1}.
    \]
    es también $R$, $f$ es holomorfa en $D(z_0, R)$ y
    \[
        f'(z) = \sum_{n=1}^{\infty} n a_n (z - z_0)^{n-1} \quad \forall z \in D(z_0, R).
    \]
\end{teorema}

\begin{proof}
    Llamemos $R'$ al radio de convergencia de la serie de las derivadas $\sum_{n=1}^{\infty} n a_n (z - z_0)^{n-1}$. Veamos en primer lugar que $R=R'$.\\
    Si $z \in D(z_0, R)$, entonces $|z-z_0| < R$. Luego $\exists z_1 \in \mathbb{C} \setminus \{z_0\}$ tal que $\sum_{n=0}^{\infty} a_n (z_1 - z_0)^n$ converge y $|z-z_0| < |z_1 - z_0| < R$. 
    \[
        0 \leq |n a_n (z - z_0)^{n}| = |n a_n \left(\frac{z - z_0}{z_1 - z_0}\right)^n (z_1 - z_0)^n| \leq M n r^2
    \]
    la última desigualdad se sigue de que $M$ es una cota superior de la sucesión $(a_n (z_1 - z_0)^n)_{n\in\mathbb{N}}$ y $r = \left|\frac{z - z_0}{z_1 - z_0}\right| < 1$.\\
    Como $\sum_{n=1}^{\infty} n r^n$ converge, por el criterio de la raíz, se tiene que $\sum_{n=1}^{\infty} M n r^n$ converge, luego por el criterio de comparación, $\sum_{n=1}^{\infty} n a_n (z - z_0)^{n-1}$ converge absolutamente. Por tanto, $R' \geq R$.\\
    Por otro lado, si $z \in \mathbb{C} \setminus \overline{D(z_0, R)}$, entonces $|z - z_0| > R$, tomando $z_1 \in \mathbb{C}$ tal que $R < |z_1 - z_0| < |z - z_0|$, sabemos que la sucesión $(a_n (z_1 - z_0)^n)_{n\in\mathbb{N}}$ no está acotada, por lo tanto $(a_{n-1} (z_1 - z_0)^{n-1})_{n\in\mathbb{N}}$ tampoco lo está. Y como
    \[
        |a_{n-1} (z_1 - z_0)^{n-1}| \leq |n a_{n-1} (z - z_0)^{n-1}| \quad \forall n \in \mathbb{N}
    \]
    deducimos que $(n a_n (z - z_0)^{n-1})_{n\in\mathbb{N}}$ no está acotada.\\
    Luego $R' \leq R$. Por tanto, $R = R'$.\\
    Para probar que $f$ es derivable en cada $z \in D(z_0, R)$, podemos suponer sin perdida de generalidad que $z_0 = 0$.\\
    Tenemos $f(z) = \sum_{n=0}^{\infty} a_n z^n$, llamemos $g(z) = \sum_{n=1}^{\infty} n a_n z^{n-1} \quad \forall z \in D(0, R)$.\\
    Sea ahora $z_1 \in D(0, R)$, nos preguntamos ¿Es $f$ derivable en $z_1$ y es $f'(z_1) = g(z_1)$? esto equivale a comprobar que:
    \[
        \lim_{z\to z_1} \frac{f(z) - f(z_1)}{z - z_1} = g(z_1).
    \]
    Tomemos $r \in \left(|z_1|, R\right)$ sea $z \in D(0, r)$:
    \[
        \left| \frac{f(z) - f(z_1)}{z - z_1} - g(z_1) \right| = \left| \frac{\sum_{n=0}^{\infty} a_n z^n - \sum_{n=0}^{\infty} a_n z_1^n}{z - z_1} - \sum_{n=1}^{\infty} n a_n z_1^{n-1} \right| \leq
    \]
    \[
        \leq \left| \frac{\sum_{n=0}^{N} a_n z^n - \sum_{n=0}^{N} a_n z_1^n}{z - z_1} + \frac{\sum_{n=N+1}^{\infty} a_n z^n - \sum_{n=N+1}^{\infty} a_n z_1^n}{z - z_1} - \sum_{n=1}^{N} n a_n z_1^{n-1} - \sum_{n=N+1}^{\infty} n a_n z_1^{n-1} \right| \leq
    \]
    \[
        \leq \left| \frac{\sum_{n=0}^{N} a_n z^n - \sum_{n=0}^{N} a_n z_1^n}{z - z_1} - \sum_{n=1}^{N} n a_n z_1^{n-1} \right| + \underbrace{\left| \sum_{n=N+1}^{\infty} a_n \frac{z^n - z_1^n}{z - z_1} \right|}_{=\sum_{n=N+1}^{\infty} a_n \left( z^{n-1} + z^{n-2}z_1 + z^{n-3}z_1^2 + \ldots + z_1^{n-1} \right)}  + \left| \sum_{n=N+1}^{\infty} n a_n z_1^{n-1} \right| \leq
    \]
    Usando aqui que $a^n-b^n = (a-b)(a^{n-1} + a^{n-2}b + \ldots + b^{n-1})$, y aplicando el binomio de Newton, tenemos que:
    \[
    \frac{z^n-z_1^n}{z-z_1} =
    \frac{(z-z_1)(z^{n-1} + z^{n-2}z_1 + z^{n-3}z_1^2 + \ldots + z_1^{n-1})}{z-z_1} =
    z^{n-1} + z^{n-2}z_1 + z^{n-3}z_1^2 + \ldots + z_1^{n-1} =
    (z-z_1)^{n-1}
    \]
    Y por tanto, sustituyendo $(z-z_1)^{n-1} = r^{n-1}$, tenemos que:
    \[
    \leq
    \left| \frac{\sum_{n=0}^{N} a_n z^n - \sum_{n=0}^{N} a_n z_1^n}{z - z_1} - \sum_{n=1}^{N} n a_n z_1^{n-1} \right| 
    + \sum_{n=N+1}^{\infty} |a_n| n r^{n-1}
    + \left| \sum_{n=N+1}^{\infty} n a_n z_1^{n-1} \right|
    \]
    Dado $\varepsilon > 0$, como $\sum_{n=1}^{\infty} n a_n z_1^{n-1}$ y $\sum_{n=1}^{\infty} |a_n| n r^{n-1}$ convergen absolutamente, $\exists N \in \mathbb{N}$ tal que:
    \[
    \sum_{n=N+1}^{\infty} |a_n| n r^{n-1} < \frac{\varepsilon}{3}, \quad \left| \sum_{n=N+1}^{\infty} n a_n z_1^{n-1} \right| < \frac{\varepsilon}{3}.
    \]
    Así, si $0 < \left| z - z_1 \right| < \delta_0 \iff z \in D(z_1, \delta_0)$, entonces $z \in D(0, r)$ y además:
    \[
        \left| \frac{f(z) - f(z_1)}{z - z_1} - g(z_1) \right| \leq \left| \frac{P_N(z) - P_N(z_1)}{z - z_1} - P_N(z_1)' \right|
        + \sum_{n=N+1}^{\infty} |a_n| n r^{n-1} + 
        \sum_{n=N+1}^{\infty} n \left| a_n \right| \left| z_1^{n-1} \right|  < \frac{\varepsilon}{3} + \frac{\varepsilon}{3} + \frac{\varepsilon}{3} = \varepsilon
    \]
    Siendo $P_N(z) = \sum_{n=0}^{N} a_n z^n$, que al ser un polinomio, es derivable en $z_1$, por lo que existe $\delta_2 > 0$ tal que si $0 < |z - z_1| < \delta_2$, entonces:
    \[
        \left| \frac{P_N(z) - P_N(z_1)}{z - z_1} - P_N(z_1)' \right| < \frac{\varepsilon}{3}.
    \]
    Asi solo basta tomar $\delta < \min\left\{ \delta_1, \delta_2 \right\}$, donde $\delta_1$ y $\delta_2$ son los $\delta$ correspondientes a las dos desigualdades anteriores, y tomar un $N$ que cumpla las dos desigualdades anteriores.
\end{proof}

Usando inducción y el teorema anterior, tenemos el siguiente teorema:

\begin{teorema}
    Sea $\sum_{n=0}^{\infty} a_n (z - z_0)^n$ una serie de potencias con radio de convergencia $R > 0$, y sea $f$ su función suma, es decir,
    \[
        f(z) = \sum_{n=0}^{\infty} a_n (z - z_0)^n \quad \forall z \in D(z_0, R).
    \]
    Entonces $f$ es infinitas veces derivable en cada $z \in D(z_0, R)$ y además:
    \[
        f^{(k)}(z) = 
        \sum_{n=0}^{\infty} \left( a_n(z-z_0)^n \right)^{(k)} =
        \sum_{n=k}^{\infty} n(n-1)(n-2)\ldots(n-k+1) a_n (z - z_0)^{n-k} \quad \forall z \in D(z_0, R) =
    \]
    \[
        = \sum_{n=k}^{\infty} \frac{n!}{(n-k)!} a_n (z - z_0)^{n-k} \quad \forall z \in D(z_0, R).
    \]
\end{teorema}

De esto último se deduce que 

\[
f^{(k)}(z_0) = k! a_k \implies a_k = \frac{f^{(k)}(z_0)}{k!} \quad \forall k \in \mathbb{N} \cup \{0\}.
\]

Veamos formas de calcular el radio de convergencia de una serie de potencias:

\begin{proposición}[Criterio de la raíz]
    Sea $\sum_{n=0}^{\infty} a_n z^n$ una serie de potencias, entonces converge absolutamente para $z$ si:
    \[
        \limsup_{n\to\infty} \sqrt[n]{|a_n z^n|} = \limsup_{n\to\infty} \sqrt[n]{|a_n|} |z| = |z| \limsup_{n\to\infty} \sqrt[n]{|a_n|} < 1.
    \]
\end{proposición}

\begin{itemize}
    \item Observemos que el radio de convergencia de $\sum_{n=0}^{\infty} a_n (z - z_0)^n$ es el mismo que el de $\sum_{n=0}^{\infty} a_n z^n$, luego podemos suponer sin pérdida de generalidad que $z_0 = 0$.
    \item La definición nos dice que
    \[
        R = \sup\left\{ |z| \geq 0 : \sum_{n=0}^{\infty} a_n z^n \text{ converge } \forall z \in D(0, r) \right\}
    \]
    \item Usando el criterio de la raiz, la serie converge absolutamente si
    \[
        |z| \overline{\lim}_{n\to\infty} \sqrt[n]{|a_n|} < 1 \iff |z| < \frac{1}{\overline{\lim}_{n\to\infty} \sqrt[n]{|a_n|}}.
    \]
    Por tanto si $\overline{\lim}_{n\to\infty} \sqrt[n]{|a_n|} \geq 0$, si $|z| < \frac{1}{\overline{\lim}_{n\to\infty} \sqrt[n]{|a_n|}}$, entonces la serie converge absolutamente.\\
    Además, si $|z| > \frac{1}{\overline{\lim}_{n\to\infty} \sqrt[n]{|a_n|}}$, entonces la sucesión $(a_n z^n)_{n\in\mathbb{N}}$ no está acotada, luego la serie $\sum_{n=0}^{\infty} a_n z^n$ diverge, con lo cual $R = \frac{1}{\overline{\lim}_{n\to\infty} \sqrt[n]{|a_n|}}$.
    \item Finalmente hemos comprobado que:
    \[
        R = 
            \begin{cases}
                \frac{1}{\overline{\lim}_{n\to\infty} \sqrt[n]{|a_n|}} & \text{ si } \overline{\lim}_{n\to\infty} \sqrt[n]{|a_n|} \in (0, \infty)\\
                0 & \text{ si } \overline{\lim}_{n\to\infty} \sqrt[n]{|a_n|} = \infty\\
                \infty & \text{ si } \overline{\lim}_{n\to\infty} \sqrt[n]{|a_n|} = 0
            \end{cases}
    \]
\end{itemize}
\begin{proposición}[Criterio del cociente]
Si $(x_n)_{n=0}^{\infty}$ es sucesión de números estrictamente positivos se cumple:

\[ \underline{\lim}_{n \to \infty} \frac{x_{n+1}}{x_n} \le \underline{\lim}_{n \to \infty} \sqrt[n]{x_n} \le \overline{\lim}_{n \to \infty} \sqrt[n]{x_n} \le \overline{\lim}_{n \to \infty} \frac{x_{n+1}}{x_n} \]

Luego si $\exists \lim_{n \to \infty} \frac{x_{n+1}}{x_n} \Rightarrow \exists \lim_{n \to \infty} \sqrt[n]{x_n}$ y $\lim_{n \to \infty} \frac{x_{n+1}}{x_n} = \lim_{n \to \infty} \sqrt[n]{x_n}$. Por lo tanto:

\[ R = \lim_{n \to \infty} \frac{|a_n|}{|a_{n+1}|} \quad \text{en caso de que exista} \]
\end{proposición}
\ejemplo{
    Halla el radio de convergencia de las siguientes series de potencias:
    \begin{itemize}
        \item Siendo la serie $\sum_{n=1}^{\infty} z^n$:\\
        Teniendo en cuenta que $a_n = 1 \quad \forall n \in \mathbb{N}$, se tiene que
        \[
            \overline{\lim}_{n\to\infty} \sqrt[n]{|a_n|} = \overline{\lim}_{n\to\infty} \sqrt[n]{1} = 1 \implies R = 1.
        \]
        \item Siendo la serie $\sum_{n=0}^{\infty} \frac{z^n}{n^n}$:\\
        Teniendo en cuenta que $a_n = \frac{1}{n^n}$, se tiene que
        \[
            \overline{\lim}_{n\to\infty} \sqrt[n]{|a_n|} = \overline{\lim}_{n\to\infty} \sqrt[n]{\frac{1}{n^n}} = \overline{\lim}_{n\to\infty} \frac{1}{n} = 0 \implies R = \infty.
        \]
        \item Siendo la serie $\sum_{n=0}^{\infty} \frac{z^2n}{n}$:\\
        Teniendo en cuenta que $a_n = \frac{1}{n}$, se tiene que
        \[
            \overline{\lim}_{n\to\infty} \sqrt[n]{|a_n|} = \overline{\lim}_{n\to\infty} \sqrt[n]{\frac{1}{n}} = 1 \implies R = 1.
        \]
    \end{itemize}
}

\begin{teorema} [Teorema de Unicidad]
    Sean $\sum_{n=0}^{\infty} a_n (z - z_0)^n$ y $\sum_{n=0}^{\infty} b_n (z - z_0)^n$ dos series de potencias que coinciden en $D(z_0, R)$ para algún $R > 0$, entonces $a_n = b_n \quad \forall n \in \mathbb{N} \cup \{0\}$.
\end{teorema}

\begin{proof}
    Por hipótesis, los radios de convergencia de ambas series son mayores o iguales que $R$. Llamemos $f$ y $g$ a las funciones suma de ambas series. Como $f \equiv g$ en $D(z_0, R)$, se tiene que $f^{(n)} \equiv g^{(n)}$ en $D(z_0, R) \quad \forall n \in \mathbb{N} \cup \{0\}$. Por tanto, 
    \[
        a_n = \frac{f^{(n)}(z_0)}{n!} = \frac{g^{(n)}(z_0)}{n!} = b_n \quad \forall n \in \mathbb{N} \cup \{0\}.
    \]
\end{proof}

